\scanpage{682}

% 682-01, upper heading
$8^\circ)$ Return to the universal case, and the homomorphism
\[
 M^{\unsure{\natural|\sharp}}\longrightarrow
 M_{\rho,\sigma}^{\unsure{\natural|\sharp}}/\mu_\rho
\]

% 682-03, upper and middle text
We now suppose that $\unsure{\mathcal G}$ corresponds to the universal case
$\mathrm{GL}(2,\mathbb Z)^{\wedge}$, with an arbitrary
$\unsure{\mathcal G}^{\unsure{\natural|\sharp}}$. We shall denote the universal
case by a subscript $0$:
$\unsure{\mathcal G}_0,\ \unsure{\mathcal G}_0^+,\
\unsure{\mathcal G}_0^{\unsure{\natural|\sharp}}$ etc., and
$\unsure{\mathcal G}',\ \unsure{\mathcal G}'^{\unsure{\natural|\sharp}}$
etc. \unsure{the|the} \unsure{former} $\unsure{\mathcal G},\
\unsure{\mathcal G}^{\unsure{\natural|\sharp}}$ etc.

% 682-04, left margin at middle
\margin{$\unsure{\swarrow}\ Z_{\rho_0,\sigma_0}$ etc $\ldots$}
\margin{We shall write $\unsure{\mathcal G}_0^+$ instead of
$\unsure{\mathcal G}_0^{+\wedge}$ $\ldots$}

% 682-06, middle-lower text and display
I set
\[
\begin{array}{rl}
(107)\qquad M^{\unsure{\natural|\sharp}}&\subset M\\[-2pt]
\Vert&\\[-6pt]
\displaystyle
\left\{u\in M\ \middle|\
\begin{array}{ll}
\text{a) }u(\rho)\text{ is }\unsure{\mathcal G}^{\unsure{\natural|\sharp}}\text{-conjugate to}
 &\varepsilon_3[\tau](\rho)\\
\text{b) }u(\sigma)\quad\rule{4.2cm}{0.3pt}&\varepsilon_2[\tau](\sigma)\\
\text{c) }u(\varepsilon_0)\quad\rule{4.2cm}{0.3pt}&\varepsilon_0^\mu\\
\text{d) }u\text{ stabilizes }\unsure{M}^{\unsure{\natural|\sharp}}&
\end{array}\right\}&
\end{array}
\]
\margin{\uline{N.B.} $M^{\unsure{\natural|\sharp}}\unsure{\supset|\ni}\mu$}

% 682-09, lower text and display
where of course
\[
\mu=\chi(u),\quad
\varepsilon_2=\varepsilon_2(u)=\varepsilon_2(\mu),\quad
\varepsilon_3=\varepsilon_3(u)=\varepsilon_3(\mu).
\]
I say that $M^{\unsure{\natural|\sharp}}$ is a \uline{subgroup} of $M$.
We can see this by \unsure{going back over} the calculations on page\ $629$ and following.
I \unsure{prefer} to proceed as follows.

\margin{681}

% Printed archive footer.
\begin{center}
681\\
\texttt{https://grothendieck.umontpellier.fr}\\
Université de Montpellier, UMR CNRS 5149, Case courrier 051, Place Eugène Bataillon\\
34095 Montpellier cedex 5 - France
\end{center}


\scanpage{683}
\mspage{644}

% 683-02, upper text and display
We have seen that we have
\[
(108)\qquad
\underset{\deleted{\unsure{\sharp}}}{Z_{\rho_0,\sigma_0}}
\hookrightarrow M[0]\inserted{{}=M(0)\cdot L_0^{\unsure{\mathcal G}}}
\simeq \widetilde M_{0,3}[0]
=\widetilde M_{0,3}(0)\cdot L_0^{\unsure{\mathcal G}_{0,3}}
\]
an inclusion which is an equality if
$\unsure{\mathcal G}_0^{\unsure{\natural|\sharp}}=\unsure{\mathcal G}_0^+$,
and which identifies $Z_{\rho,\sigma}$ with
\[
M^{!}[0]\simeq M(0)\cdot L_0^\pi
\quad\text{if}\quad
\unsure{\mathcal G}_0^{\unsure{\natural|\sharp}}=\pi\unsure{:=}\Pi_0.
\]

% 683-08, bracketed passage
[ --- in fact, if
$\unsure{\mathcal G}_0^{\unsure{\natural|\sharp}}\subset\Pi_0$,
then we find
\[
\begin{split}
(109)\quad Z_{\rho,\sigma}=\left\{
u=\underset{\substack{\rotatebox{-90}{$\in$}\\M(0)}}{u_0}\cdot
\underset{\substack{\rotatebox{-90}{$\in$}\\\deleted{\unsure{\mathcal G}_0}L_0^\pi}}{\ell}
\ \middle|\
\begin{array}{l}
\text{a) }u(\rho)=\alpha(\rho)\quad\text{with }\alpha\in
\unsure{\mathcal G}_{\unsure{\bar\tau}}^{\unsure{\natural|\sharp}}\\
\text{b) }u(\sigma)=\beta(\sigma)\quad\text{with }\beta\in
\unsure{\mathcal G}_{\unsure{\bar\tau}}^{\unsure{\natural|\sharp}}\\
\text{d) }u\unsure{\varepsilon|\in}\text{ stabilizes }
\unsure{\mathcal G}_0^{\unsure{\natural|\sharp}}
\end{array}\right\}
\end{split}
\]
\margin{this explicit description does not \unsure{seem} to be of any use.}

% 683-12, middle-lower text with cancellations and additions
and we see that $\alpha,\beta$ in a) b) are determined by
$u=u_0\ell$ uniquely
(\deleted{as soon as} \unsure{even}
\unsure{if we|we} suppose only $u\in M^{!}[0]$ and no more,
and \inserted{\unsure{that we require}}
\deleted{\unsure{$\alpha\in$|with} $\unsure{\mathcal G}_0^{\unsure{\natural}}$}
$\alpha,\unsure{\beta}\in\unsure{\Pi_{0\bar\tau}}$,
\deleted{\unsure{can|In short}}
We know that this is a \unsure{subgroup|subgroup}
\inserted{of $\operatorname{Aut}(\unsure{\mathcal G}_0^+)\times
\widehat{\mathbb Z}^{*}\times\mu\times\mu$
--- and in fact, of $\operatorname{Aut}(\unsure{\mathcal G}_0^+)$}
(proposition\ page\ $629$).
\deleted{\unsure{Let us consider}}

% 683-19, lower paragraph
As this group operates on \deleted{$M$ and}
$\unsure{\mathcal G}$ via $\widetilde M_{0,3}$,
\inserted{$=\text{ via }\widetilde M_{0,3}^{!}[0]$}
it maps into $\widetilde M_{0,3}^{!}[0]$ by a homomorphism of groups,
hence into \deleted{$\unsure{\widetilde M_{0,3}[0]}$}
$M[0]^{!}\simeq\widetilde M_{0,3}[0]$ by a homomorphism of groups. ]

\margin{682}

% Printed archive footer.
\begin{center}
682\\
\texttt{https://grothendieck.umontpellier.fr}\\
Université de Montpellier, UMR CNRS 5149, Case courrier 051, Place Eugène Bataillon\\
34095 Montpellier cedex 5 - France
\end{center}


\scanpage{684}

% 684-01, upper text and tuple
\deleted{\unsure{The homomorphism}} Let us consider the system
\[
\begin{split}
(110)\quad\bigl(&\inserted{\unsure{\mathcal G}_0^{+\wedge},}
Z_{\rho_0,\sigma_0},\quad
L_0^{\unsure{\natural|\sharp}}\subset Z_{\rho_0,\sigma_0},\quad
Z_{\rho_0,\sigma_0}\longrightarrow
\operatorname{Aut}(\unsure{\mathcal G}_0^{\unsure{+\wedge}}),\quad
L_0^{\unsure{\natural|\sharp}}\hookrightarrow\unsure{\mathcal G}_0^{+\wedge}\bigr)
\end{split}
\]
(cf. ``digression'' below $9^\circ$) which gives rise to
\[
(110')\qquad
G_{\rho_0,\sigma_0},\quad
\unsure{\mathcal G}_0^{+\wedge}\subset G_{\rho_0,\sigma_0},\quad
Z_{\rho_0,\sigma_0}\hookrightarrow G_{\rho_0,\sigma_0},
\]
and let us consider the system analogous \inserted{to $(110^{*})$}
\[
\begin{split}
(111^{\unsure{*|\sharp}})\quad\bigl(&\inserted{\unsure{\mathcal G}_0^+,}
\underset{\substack{\Vert\\M(0)\cdot L_0^\pi}}{M[0]^{\unsure{!|\natural}}},\quad
\underset{\text{(\unsure{generated} by }\varepsilon_0\text{)}}{L_0^{\unsure{\mathcal G}}}
\subset M[0]^{\unsure{!|\natural}},\quad
\unsure{\deleted{M[0]}}^{\unsure{!|\natural}}\longrightarrow
\operatorname{Aut}(\unsure{\mathcal G}_0^+),\quad
L_0^{\unsure{\mathcal G}}\hookrightarrow\unsure{\mathcal G}_0^+\bigr)
\end{split}
\]
which gives rise (loc.\ cit ($9^\circ$)) to the situation
\[
(111')\qquad
\bigl(M,\quad\unsure{\mathcal G}_0^+\subset M,\quad
M[0]'\hookrightarrow\unsure{\mathcal G}_0^+\bigr).
\]

% 684-14, lower paragraph
Now thanks to $(108)$, we have a homomorphism from $(\unsure{110|109|119})$ into $(111)$,
corresponding to the identity on $\unsure{\mathcal G}_0^+$,
$(108)$ on $Z_{\rho_0,\sigma_0}$, the canonical inclusion\ :
\deleted{\unsure{on}}
$L_0^{\unsure{\natural|\sharp}}\subset L_0^{\unsure{\mathcal G}}$.
This yields a homomorphism from $(110')$ into $(111')$, in particular a canonical
homomorphism of exact sequences

\margin{683}

% Printed archive footer.
\begin{center}
683\\
\texttt{https://grothendieck.umontpellier.fr}\\
Université de Montpellier, UMR CNRS 5149, Case courrier 051, Place Eugène Bataillon\\
34095 Montpellier cedex 5 - France
\end{center}


\scanpage{685}
\mspage{645}

% 685-02, upper commutative diagram
% (112): the first \mathcal G_0^+ seems written over an earlier G_\rho (unresolved).
\begin{samepage}
(112)\quad
\diagram{\begin{tikzcd}[column sep=2.1em,row sep=2.6em]
1 \arrow[r,no head] &
\unsure{\mathcal G}_0^+ \arrow[r] \arrow[d,"\unsure{1}"] &
G_{\rho_0,\sigma_0} \arrow[r] \arrow[d,hook] &
\gamma_{\rho_0,\sigma_0} \arrow[r] \arrow[d,hook,dotted] \arrow[d] & 1 \\
1 \arrow[r,no head] &
\unsure{\mathcal G}_0^+ \arrow[r,no head] & M \arrow[r] &
\widetilde\Gamma_{\mathbb Q} \arrow[r,no head] & 1
\end{tikzcd}}
\end{samepage}

% 685-03, explanatory diagram and text
where
\diagram{\begin{tikzcd}[column sep=4em,row sep=1em]
\gamma_{\rho_0,\sigma_0} \arrow[r] \arrow[d,phantom,"\unsure{\wr}" description] &
\widetilde\Gamma_{\mathbb Q} \arrow[d,phantom,"\unsure{\wr}" description] \\
Z_{\rho_0,\sigma_0}/L_0^{\unsure{\mu}} &
M[0]/L_0^{\unsure{\mathcal G}}
\end{tikzcd}}
is \inserted{\unsure{the one}} induced by
$(\unsure{108})$, and is identified with:
\[
(113)\qquad Z_{\rho_0,\sigma_0}/L_0^{\unsure{\natural|\sharp}}
\longrightarrow M[0]/L_0^{\unsure{\mathcal G}}
\]

% 685-07, middle text
By the very definition of
$L_0^{\unsure{\natural|\sharp}}$ as being
$Z_{\rho_0,\sigma_0}\cap\unsure{\mathcal G}_0^+$
\deleted{\unsure{and|this}}, we see that $(113)$ is injective\ hence
$G_{\rho_0,\sigma_0}\longrightarrow M$ is always injective
(it is an isomorphism\ if
$\unsure{\mathcal G}_0^{\unsure{\natural|\sharp}}\supset\Pi_0$
$\ldots$).
We deduce by comparison from $(112)$ a homomorphism of exact sequences

% 685-09, lower commutative diagram
\begin{samepage}
(114)\quad
\diagram{\begin{tikzcd}[column sep=2.1em,row sep=2.6em]
1 \arrow[r,no head] &
\unsure{\mathcal G}_0^{\unsure{\natural|\sharp}} \arrow[r] \arrow[d,hook] &
G_{\rho_0,\sigma_0}^{\unsure{\natural|\sharp}} \arrow[r] \arrow[d,hook] &
\gamma_{\rho_0,\sigma_0} \arrow[r,no head] \arrow[d,hook] & 1 \\
1 \arrow[r,no head] &
\unsure{\mathcal G}_0^{\unsure{+}} \arrow[r] &
M \arrow[r] &
\deleted{\unsure{\gamma}}\widetilde\Gamma_{\mathbb Q} \arrow[r,no head] & 1.
\end{tikzcd}}
\end{samepage}

% 685-10, bottom text
We then see immediately that the image of
$G_{\rho_0,\sigma_0}^{\unsure{\natural|\sharp}}$ in $M$ is none other than
$M^{\unsure{\natural|\sharp}}$ $(107)$. We therefore have

\margin{684}

% Printed archive footer.
\begin{center}
684\\
\texttt{https://grothendieck.umontpellier.fr}\\
Université de Montpellier, UMR CNRS 5149, Case courrier 051, Place Eugène Bataillon\\
34095 Montpellier cedex 5 - France
\end{center}


\scanpage{686}

\[
(115)\qquad
\left\{
\begin{array}{ll}
\underset{\substack{\text{\unsure{defined in}}\\(107)}}{M^{\unsure{\natural|\sharp}}}
  \simeq G_{\rho_0,\sigma_0}^{\unsure{\natural|\sharp}},\\[5pt]
\gamma_{\rho_0,\sigma_0}\hookrightarrow\widetilde{\Gamma}_{\mathbb Q}
  \deleted{\unsure{\Leftarrow}\{u\in\widetilde{\Gamma}_{\mathbb Q}}
\end{array}
\right.
\]

\margin{
\diagram{
\begin{tikzcd}[row sep=small]
\text{\uline{to be checked}}\arrow[d]\\
\text{\uline{N.B.}\quad We have}
\end{tikzcd}
}
\[
(116)\quad
\left\{
\begin{array}{l}
\gamma_{\rho_0,\sigma_0}
 \deleted{\unsure{=\Gamma_0}}
 \inserted{\unsure{=\Gamma_0^{\natural}|\Gamma_0^{\sharp}}}
 =\left\{u\in\widetilde{\Gamma}_{\mathbb Q}\ \middle|\
 \begin{array}{l}
\text{the lifting }u_0\\
\text{of }u\text{ to }u_0\in M(0)\\
\text{satisfies the}\\
\text{conditions a) b) d)}\\
\text{of (107)}
\end{array}\right\}\\
\text{(\unsure{notations}}\\
\gamma_0=\Gamma_0^{\unsure{\natural|\sharp}},
\ \text{where }\Gamma_0\text{ denotes }\widetilde{\Gamma}_{\mathbb Q}\text{)}
\end{array}
\right.
\]
}

which \unsure{shows} at the same time in the
\unsure{process} that $M^{\unsure{\natural|\sharp}}$ is indeed a
subgroup of $M$.

We now have
\[
(\unsure{117})\qquad
M_{\rho_0,\sigma_0}^{\unsure{\natural|\sharp}}
\overset{\unsure{\mathrm{dfn}|\mathrm{déf}}}{\simeq}
\unsure{N|M}_{\rho_0}^{\unsure{\natural|\sharp}}
\times_{\gamma_{\rho_0,\sigma_0}}
N_{\sigma_0}^{\unsure{\natural|\sharp}}
\times_{\unsure{\gamma_{\rho_0,\sigma_0}}}
\underbrace{G_{\rho_0,\sigma_0}^{\unsure{\natural|\sharp}}}
 _{\simeq M^{\unsure{\natural|\sharp}}},
\]
where $G_{\rho_0,\sigma_0}^{\unsure{\natural|\sharp}}\simeq
M^{\unsure{\natural|\sharp}}$ \unsure{therefore|gives}
\[
\begin{aligned}
N_{\rho_0}^{\unsure{\natural|\sharp}}
 &\unsure{\simeq|=}\left\{u\in M^{\unsure{\natural|\sharp}}\mid
 u(\rho)=\rho^{\varepsilon_3(u)}\right\}\\
N_{\sigma_0}^{\unsure{\natural|\sharp}}
 &\unsure{\simeq|=}\left\{u\in M^{\unsure{\natural|\sharp}}\mid
 u(\sigma)\unsure{\equiv|=}\sigma^{\varepsilon_2(u)}\right\}
\end{aligned}
\]
and finally we have
\[
(118)\qquad
\left\{
\begin{array}{l}
M_{\rho_0,\sigma_0}^{\unsure{\natural|\sharp}}
\simeq
N_{M^{\unsure{\natural|\sharp}}}(\rho)
\times_{\deleted{\unsure{\gamma_{\rho_0,\sigma_0}}}
 \inserted{\Gamma_0^{\unsure{\natural|\sharp}}}}
N_{M^{\unsure{\natural|\sharp}}}(\sigma)
\times_{\deleted{\unsure{\gamma_{\rho_0,\sigma_0}}}
 \inserted{\Gamma_0^{\unsure{\natural|\sharp}}}}
M^{\unsure{\natural|\sharp}}.\\[5pt]
\text{where }\Gamma_0^{\unsure{\natural|\sharp}}\unsure{\ne|=}\gamma_0
\simeq\gamma_{\rho_0,\sigma_0}\subset\widetilde{\Gamma}_{\mathbb Q}
\quad\text{\unsure{made explicit} in (116)},
\end{array}
\right.
\]
\uuline{whence} an \deleted{isomorphism of} \unsure{exact|exact}
\unsure{sequences|sequence}
\begin{samepage}
\[
(119)
\]
\diagram{
\begin{tikzcd}[column sep=small]
1\arrow[r]&
\overbrace{SZ_{\rho_0}\times SZ_{\sigma_0}}^{\inserted{L_\rho\times L_\sigma}}
\arrow[r]&
M_{\rho_0,\sigma_0}^{\unsure{\natural|\sharp}}
\arrow[r]\arrow[d,"\inserted{\unsure{\natural|\sharp}}"']&
\deleted{\unsure{\gamma}}M^{\unsure{\natural|\sharp}}\arrow[r]&1\\
&&{}&&
\end{tikzcd}
}
\end{samepage}
where the kernel $SZ_{\rho_0}\times SZ_{\sigma_0}$ is a
\unsure{very small} group, namely

\margin{685}

% Printed archive footer.
\begin{center}
685\\
\texttt{https://grothendieck.umontpellier.fr}\\
Université de Montpellier, UMR CNRS 5149, Case courrier 051, Place Eugène Bataillon\\
34095 Montpellier cedex 5 - France
\end{center}


\scanpage{687}
\mspage{646}

\[
(120)\quad
\operatorname{Ker}\left(
M_{\rho_0,\sigma_0}^{\unsure{\natural|\sharp}}
\overset{\text{epimorphism}}{\longrightarrow}M^{\unsure{\natural|\sharp}}
\right)
\simeq SZ_{\rho_0}\times SZ_{\sigma_0}
\simeq L_\rho\times L_\sigma
\simeq \mathbb Z/6\mathbb Z\times\mathbb Z/4\mathbb Z .
\]

We are going to improve \unsure{further|further} the \deleted{\illegible{2}}\inserted{situation},
by cutting out of $M_{\rho_0,\sigma_0}^{\unsure{\natural|\sharp}}$
a group \unsure{even|even} smaller, which is an extension of
$M^{\unsure{\natural|\sharp}}$ by $\mu$.
\deleted{For this, we define}
\[
\deleted{
N_{\unsure{\sigma_0|\rho_0}}^{\unsure{\natural|\sharp}!}
=N_{\sigma}^{\unsure{\natural|\sharp}}\cap M_0^!
\unsure{=|\cdot}M^{\unsure{\natural|\sharp}}\unsure{\cap}
}
\]

\unsure{(}

Let us note that
\[
(121)\qquad
\begin{array}{c}
M_{\rho_0,\sigma_0}^{\unsure{\natural|\sharp}}
=N_{\rho_0}^{\unsure{\natural|\sharp}}
\times_{\Gamma_0^{\unsure{\natural|\sharp}}}
N_{\sigma_0}^{\unsure{\natural|\sharp}}
\times_{\Gamma_0^{\unsure{\natural|\sharp}}}
\deleted{\unsure{G_{\rho_0,\sigma_0}^{\natural}|G_{\rho_0,\sigma_0}^{\sharp}}}
\underbrace{\inserted{\unsure{G|M}_{\rho_0,\sigma_0}^{\unsure{\natural|\sharp}}}}
 _{M^{\unsure{\natural|\sharp}}}\\[5pt]
\cap\\[3pt]
N_{\rho_0}\times_{\Gamma_0}N_{\sigma_0}\times_{\Gamma_0}M .
\end{array}
\]
\margin{where $\Gamma_0\overset{\unsure{\mathrm{dfn}|\mathrm{déf}}}{=}
\widetilde{\Gamma}_{\mathbb Q}$, \unsure{and}}

\inserted{and}\ \unsure{where} this last fiber product is none other than the
quantity $M_{\rho_0,\sigma_0}^{(\unsure{\natural|\sharp}\,')}$
which would correspond \unsure{to the|in the} case
\deleted{$\unsure{M_{\rho_0,\sigma_0}}$}
of $\unsure{\mathcal G}_0^{\unsure{\natural|\sharp}\,'}
=\unsure{\mathcal G}_0^+$ in \unsure{all|all} its entirety
(denoted $M_{\rho,\sigma}$ on \unsure{page|page} 626 and following).
In this group we defined the subgroup, which was denoted
$M_{\rho,\sigma}^{\unsure{\natural|\sharp}}$
\deleted{\unsure{of}}\inserted{\unsure{on}} page\ $626'$,
and which I am now going to denote instead by
$M_{\rho_0,\sigma_0}^{\unsure{?|!}}$, defined as

\margin{686}

% Printed archive footer.
\begin{center}
686\\
\texttt{https://grothendieck.umontpellier.fr}\\
Université de Montpellier, UMR CNRS 5149, Case courrier 051, Place Eugène Bataillon\\
34095 Montpellier cedex 5 - France
\end{center}


\scanpage{688}

\[
\deleted{M_{\unsure{\rho_0|\rho_0,\sigma_0}}^{\unsure{\natural|\sharp}}}
\qquad
\deleted{M_{\rho_0,\sigma_0}^{?}=\unsure{Z_{\rho_0}^{!}}}
\]

\[
(122)\qquad
\deleted{M_{\unsure{\rho_0|\rho_0,\sigma_0}}}
M^?=N_{\unsure{\rho|\rho_0}}^?
\times_{\Gamma_0}Z_{\sigma_0}^!\times_{\Gamma_0}M
\]
where, we recall that

\diagram{
\begin{tikzcd}[column sep=small,row sep=small]
&
N_{\rho_0}=\operatorname{Norm}_M(\unsure{L}_{\rho_0}),\quad
N_{\unsure{\sigma|\sigma_0}}=\operatorname{Norm}_M(L_{\sigma_0})
\\
(123)\quad\Biggl\{&
\deleted{\unsure{Z_{\rho_0}}}\qquad
Z_{\unsure{\sigma|\sigma_0}}
=\operatorname{Centr}_M(L_{\sigma_0})
=\operatorname{Centr}_M(\sigma_0)
\\
&
\deleted{N}\inserted{\unsure{Z}}_{\sigma_0}^!
=Z_{\sigma_0}\cap M_0^!
\left(=N_{\sigma_0}\cap M_0^!\right)\subset N_{\sigma_0}
\\
&
N_{\rho_0}^?=Z_{\rho_0}^!\cdot\mu_{\unsure{\tau_1|\tau'}}
\subset N_{\rho_0}\unsure{)|,}\quad
\text{where }Z_{\rho_0}^!=Z_{\rho_0}\cap M^!\ )
\arrow[uu,bend left=45,start anchor=west,end anchor=west]
\end{tikzcd}
}

so that we do indeed have
\[
M^?\subset N_{\rho_0}\times_{\gamma_0}N_{\sigma_0}
\times_{\Gamma_0}M .
\]
\margin{\unsure{independent|independently} of the choice of a
$\unsure{\mathcal G}^{\unsure{\natural|\sharp}}$}

This recalled, we are going to set
\begin{samepage}
\[
(124)
\]
\diagram{
\begin{tikzcd}[column sep=small,row sep=small,nodes in empty cells]
M_{\rho_0,\sigma_0}^{\unsure{\natural|\sharp}
 \deleted{\unsure{!|\natural}}?}
\arrow[r,phantom,"\subset" description]\arrow[d,equal]&
M_{\rho_0,\sigma_0}^{\unsure{\natural|\sharp}}
\\
M_{\rho_0,\sigma_0}^{\unsure{\natural|\sharp}}\cap M^?\arrow[d,equal]&
\begin{gathered}
\text{(\unsure{intersection}}\\
\text{in }N_{\rho_0}\times_{\Gamma_0}N_{\sigma_0}
\times_{\Gamma_0}M\text{)}
\end{gathered}
\\
\deleted{\left(
\unsure{N|M}_{\rho_0}^{\unsure{!|?}}
\times_{\Gamma_0}N_{\sigma_0}^{!}
\times_{\Gamma_0}M^{\unsure{\natural|\sharp}}
\right)}
&
\end{tikzcd}
}
\end{samepage}
and let us consider \unsure{the homomorphism|the homomorphism} induced by (119)
\begin{samepage}
\[
(125)
\]
\diagram{
\begin{tikzcd}
M_{\rho_0,\sigma_0}^{\unsure{\natural|\sharp}?}
\arrow[r]\arrow[d,hook]&
M^{\unsure{\natural|\sharp}}\arrow[d,hook]\\
M^?\arrow[r]&M
\end{tikzcd}
}
\end{samepage}
.

\margin{687}

% Printed archive footer.
\begin{center}
687\\
\texttt{https://grothendieck.umontpellier.fr}\\
Université de Montpellier, UMR CNRS 5149, Case courrier 051, Place Eugène Bataillon\\
34095 Montpellier cedex 5 - France
\end{center}


\scanpage{689}
\mspage{647}

We have seen (page\ $626'$, (55)) that $M^?\longrightarrow M$ makes $M^?$
an extension of $M$ with kernel $\mu_\rho$.
Thus the kernel of $M_{\rho_0,\sigma_0}^?\longrightarrow
M^{\unsure{\natural|\sharp}}$ is equal to
$\mu_\rho\cap M_{\rho_0,\sigma_0}^?$
now it is clear that we have
$\mu_\rho\subset M_{\rho_0,\sigma_0}^{\unsure{\natural|\sharp}}$
\inserted{cf. (\unsure{119|114})},
hence $\mu_\rho\subset
M_{\rho_0,\sigma_0}^{\unsure{\natural|\sharp}?}$.

We thus find a \deleted{\unsure{exact sequence}}\inserted{homomorphism}
of exact sequences
\begin{samepage}
\[
(126)
\]
\diagram{
\begin{tikzcd}[column sep=small,row sep=small]
&&\inserted{M_{\rho_0,\sigma_0}^{\unsure{\natural|\sharp}}}
\arrow[d,phantom,"\cup" description]&&\\
1\arrow[r]&\mu_\rho\arrow[r]\arrow[d,equal]&
M_{\rho_0,\sigma_0}^{\inserted{\unsure{\natural|\sharp}}?}
\arrow[r]\arrow[d,hook]&
M^{\unsure{\natural|\sharp}}\arrow[d,hook]&\\
1\arrow[r]&\mu_\rho\arrow[r]&M^?\arrow[r]&M\arrow[r]&1
\end{tikzcd}
}
\end{samepage}
and it remains to see whether \unsure{the} \deleted{\unsure{arrow}}\inserted{homomorphism}
\[
(\infty)\qquad
M_{\rho_0,\sigma_0}^?\twoheadrightarrow M^{\unsure{\natural|\sharp}}
\]
is an epimorphism, and if not, \unsure{characterize} its image.
\deleted{\unsure{Consid|Let us consider}}
Let us note that \inserted{the definition\unsure{,} (124) also gives}
\deleted{\unsure{this} \unsure{homomorphism} \unsure{fits} into a
\unsure{diagram}}
taking into account the expressions (118) for
\deleted{\unsure{of|some} exact sequences}
$M_{\rho_0,\sigma_0}^{\unsure{\natural|\sharp}}$ and (122) for $M^?$ :
\[
(127)\qquad
M_{\rho_0,\sigma_0}^{\unsure{\natural|\sharp}?}
\simeq N_{\rho_0}^{\unsure{\natural|\sharp}?}
\times_{\deleted{\unsure{\gamma_{\rho_0,\sigma_0}}}
 \inserted{\Gamma_0^{\unsure{\natural|\sharp}}}}
\deleted{N}\inserted{\unsure{Z}}_{\sigma_0\,0}^{\unsure{\natural|\sharp}!}
\times_{\deleted{\unsure{\gamma_{\rho_0,\sigma_0}}}
 \inserted{\Gamma_0^{\unsure{\natural|\sharp}}}}
M^{\unsure{\natural|\sharp}}
\]
\margin{
\diagram{
\begin{tikzcd}[column sep=small,row sep=small]
\gamma_0\arrow[r,phantom,"\subset" description]\arrow[d,equal]&
\widetilde{\Gamma}_{\mathbb Q}\\
\operatorname{Im}(M^{\unsure{\natural|\sharp}})&
\end{tikzcd}
}
}
where we have set
\[
(128)\qquad
\left\{
\begin{aligned}
N_{\rho_0}^{\unsure{\natural|\sharp}?}
&=N_{\rho_0}^{?}\cap N_{\rho_0}^{\unsure{\natural|\sharp}}
=N_{\rho_0}^{\unsure{?|!}}\cap M^{\unsure{\natural|\sharp}}\\
\deleted{N}\inserted{\unsure{Z}}_{\sigma_0\,0}^{\unsure{\natural|\sharp}!}
&=Z_{\sigma_0\,0}^!\cap N_{\sigma_0}^{\unsure{\natural|\sharp}}
=\deleted{N}\inserted{\unsure{Z}}
_{\sigma_0\,\underset{\deleted{\unsure{?|\natural}}}{0}}^!
\cap M^{\unsure{\natural|\sharp}}
\end{aligned}
\right.
\]

\margin{688}

% Printed archive footer.
\begin{center}
688\\
\texttt{https://grothendieck.umontpellier.fr}\\
Université de Montpellier, UMR CNRS 5149, Case courrier 051, Place Eugène Bataillon\\
34095 Montpellier cedex 5 - France
\end{center}


\scanpage{690}

and the homomorphism $(\unsure{\infty|\alpha})$ is none other than the third projection of the fiber product.
Thus the image of $(\unsure{\infty|\alpha})$ is none other than the inverse image, in $M^{\unsure{\natural|\sharp}}$ under $M^{\unsure{\natural|\sharp}}\longrightarrow\gamma_0=\gamma_{\deleted{\unsure{\natural|\sharp}}\rho_0,\sigma_0}$ of the image of
\[
 (\unsure{\infty\infty|\alpha\alpha})\qquad
 N_{\rho_0}^{\deleted{\unsure{\natural|\sharp}}\inserted{\unsure{\natural|\sharp}?}}
 \mathop{\times}_{\deleted{\unsure{\gamma_{\rho_0,\sigma_0}}}\inserted{\Gamma_0^{\unsure{\natural|\sharp}}}}
 N_{\sigma_0}^{\unsure{\natural|\sharp}}{}_{0}^{!}
 \longrightarrow
 \deleted{\unsure{\gamma_{\rho_0,\sigma_0}}}\inserted{\Gamma_0^{\unsure{\natural|\sharp}}}
 (\subset\Gamma_0)
\]
It thus remains to show that this homomorphism is surjective, or again that the homomorphisms
\diagram{\(\left\{\vcenter{\hbox{\begin{tikzcd}[row sep=large,column sep=huge]
 N_{\rho_0}^{\unsure{\natural|\sharp}?} \arrow[r] & \deleted{\unsure{\gamma_{\rho_0,\sigma_0}}}\inserted{\Gamma_0^{\unsure{\natural|\sharp}}} \\
 N_{\sigma_0}^{\unsure{\natural|\sharp}}{}_{0}^{!} \arrow[r] & \Gamma_0^{\unsure{\natural|\sharp}}
\end{tikzcd}}}\right.\)}
are surjective (the first will then be an extension of $\Gamma_0^{\unsure{\natural|\sharp}}$ by $\mu_\rho$, the second an isomorphism). So let $u\in\Gamma_0^{\unsure{\natural|\sharp}}$, \deleted{defined by a}\inserted{coming from} $u_{\unsure{0}}\unsure{\in}M^{\unsure{\natural|\sharp}}\deleted{\unsure{_{\rho,\sigma}}}\simeq G_{\rho,\sigma}^{\unsure{\natural|\sharp}}$ thus satisfying conditions a) b) c) d) of (107). We thus see that there exists $\beta_0\in\unsure{\mathcal G}_0^{\unsure{\natural|\sharp}}\inserted{\subset M^{\unsure{\natural|\sharp}}}$ such that $\beta^{-1}u_0\in N(\sigma_0)$.

\margin{Here I have avoided the verifications, thinking I would return to them in the following pages, it turns out that this is not urgent -- it was more important to develop the special cases ``Fermat \unsure{tower|towers}'' to clarify the ideas \inserted{of the \unsure{general} formulations of the $M_{\rho,\sigma}$\ldots.}}

\margin{689}

% Printed archive footer.
\begin{center}
689\\
\texttt{https://grothendieck.umontpellier.fr}\\
Université de Montpellier, UMR CNRS 5149, Case courrier 051, Place Eugène Bataillon\\
34095 Montpellier cedex 5 - France
\end{center}


\scanpage{691}
\mspage{648}

$9^{\circ})$ \uline{Digression on group formalism}.

\textcircled{a}

Let
\[
 (1)\qquad
 \text{\uline{$G$}},\quad
 \text{\uline{$\unsure{\mathcal G}\subset G$}},\quad
 \text{\uline{$\unsure{N}\overset{\varphi_N}{\hookrightarrow}G$}},\quad
 \inserted{\text{hence}}\quad
 SN=\deleted{\unsure{N\cap\mathcal G}}\inserted{\varphi^{-1}(\unsure{\mathcal G})}\subset N
\]
be the situation of a group $G$, invariant subgroup $\unsure{\mathcal G}$, and subgroup $N$ such that
\[
 (\unsure{\alpha})\qquad
 N\xrightarrow{\text{epimorphism}}G/\unsure{\mathcal G}\deleted{\unsure{\overset{\mathrm{dfn}}{=}\gamma}},
 \qquad\text{i.e.}\quad N/SN\xrightarrow{\sim}G/\unsure{\mathcal G}
\]
\deleted{\unsure{is an epimorphism}} We are going to reconstruct this situation from the following one
\[
 (2)\qquad
 \unsure{\mathcal G},\quad N,\quad
 SN\subset N\deleted{\unsure{\operatorname{Aut}}},\quad
 N\xrightarrow{\deleted{\unsure{\varphi}}\inserted{\Theta}}\operatorname{Aut}(\unsure{\mathcal G}),
 \quad\inserted{SN\xrightarrow[\varphi=\varphi_{SN}]{\deleted{\unsure{\varphi}}}\unsure{\mathcal G}}
\]
where $\unsure{\mathcal G},N$ are groups, $SN$ a normal subgroup of $N$, $N\longrightarrow\operatorname{Aut}(\unsure{\mathcal G})$ a homomorphism of groups, i.e. an \uuline{action} of $N$ on $\unsure{\mathcal G}$; finally $SN\longrightarrow\unsure{\mathcal G}$ a homomorphism of groups. We assume that \deleted{the action}
\deleted{that $SN\longrightarrow\unsure{\mathcal G}$ commutes}
\[
 (3)\quad
 \left\{\begin{aligned}
 &\inserted{\text{a) the action}}\text{ of }SN\text{ on }\unsure{\mathcal G}\text{ \unsure{induced|induced} by that of }\unsure{\mathcal G|G},\\
 &\text{is also that deduced from }SN\longrightarrow\unsure{\mathcal G},\\
 &\text{via the adjoint operation of }\unsure{\mathcal G},\\
 &\inserted{\text{b) }SN\longrightarrow\unsure{\mathcal G}\text{ is compatible with the actions of }N.}
 \end{aligned}\right.
\]
and \deleted{\unsure{It is|I see}} \deleted{\unsure{It is|I see}} clear how (1) defines (2), let us show how a system (2) allows one to reconstruct a system (1).

\margin{690}

% Printed archive footer.
\begin{center}
690\\
\texttt{https://grothendieck.umontpellier.fr}\\
Université de Montpellier, UMR CNRS 5149, Case courrier 051, Place Eugène Bataillon\\
34095 Montpellier cedex 5 - France
\end{center}


\scanpage{692}

We thus start from \unsure{(2)}, we set
\[
 (4)\qquad S_0G=N.\unsure{\mathcal G}\qquad(\text{\unsure{Ch direct|direct product}}),
\]
we define a homomorphism
\diagram{\((5)\quad\left\{\vcenter{\hbox{\begin{tikzcd}[row sep=large,column sep=large]
 SN \arrow[r,hook,"i"] & S_0G \\
 \unsure{u} \arrow[r,mapsto] & \unsure{u}.\varphi(\unsure{u})^{-1}=\varphi(u)^{-1}.u
\end{tikzcd}}}\right.\)}

\margin{\uline{N.B.} we have \inserted{(by hypothesis (a))} for $u\in SN$, $x\in\unsure{\mathcal G}$
\[
 (*)\quad\left\{u(x)=\underset{\substack{\text{in }S_0G\\=N.\unsure{\mathcal G}}}{\underbrace{uxu^{-1}}}=\varphi(u)x\varphi(u)^{-1}\right.
\]
so if $x=\varphi(u)$, we have
\[
 u(x)=uxu^{-1}=x
\]
i.e. $x=\varphi(u)$ and $u$ commute.}

Let us check that it is a homomorphism
\deleted{\unsure{we}}
\[
 i(\unsure{u'}u)=(u'u)\varphi(u'u)^{-1}
\]
\[
\begin{aligned}
 i(u')i(u)&=u'\varphi(u')^{-1}u\varphi(u)^{-1}\\
 &\deleted{=u'\varphi(u')^{-1}u\varphi(u')}\\
 &=u'u\underset{\substack{=\varphi(u^{-1})\varphi(u'^{-1})\varphi(u)\\\text{by a)}}}{\underbrace{(u^{-1}\varphi(u'^{\unsure{-1}})u)}}\varphi(u^{-1})^{\deleted{\unsure{-1}}}\\
 &=u'u\varphi(u^{-1}u'^{-1}\deleted{uu^{-1}})=(u'u)\varphi((u'u)^{-1})=(u'u)\varphi(u'u)^{-1}\qquad\text{ok.}
\end{aligned}
\]

We have not yet used the fact that is normal in $N$, \inserted{and hypothesis (b)} this implies that the image of \unsure{(4)|(5)} is normal in $S_0G$, for \deleted{it is} we have, if $a\in N$, $u\in SN$
\[
 \deleted{\unsure{(6)(6)}}\qquad\deleted{\unsure{(*)}}\quad
 a(u.\varphi(u)^{-1})a^{-1}
 =\underbrace{aua^{-1}}_{\in SN}.\underset{=\varphi(au^{-1}a^{-1})=\varphi(au^{-1}a^{-1})^{-1}}{\underbrace{a\varphi(u^{-1})a^{-1}}}
\]
i.e. the homomorphism (5) \deleted{\unsure{is}} commutes with the actions of $N$ on $SN$ and on $S_0G\supset N$ -- thus the image of $SN$ is normalized by $N\subset S_0G$. We still have to check that it is also normalized by $\unsure{\mathcal G}$, \deleted{so let $x\in\unsure{\mathcal G}$, we have}\inserted{but we have better:}
\[
 \deleted{xu\varphi(u^{-1})x^{-1}=(xux^{-1})(x\varphi(u^{-1})x^{-1}}
\]
The image of $SN\xrightarrow{i}S_0G$ \uline{commutes} with $\unsure{\mathcal G}$, i.e.

\margin{691}

% Printed archive footer.
\begin{center}
691\\
\texttt{https://grothendieck.umontpellier.fr}\\
Université de Montpellier, UMR CNRS 5149, Case courrier 051, Place Eugène Bataillon\\
34095 Montpellier cedex 5 - France
\end{center}


\scanpage{693}
\mspage{649}

$\forall u\in SN$, $u\varphi(u)^{-1}$ induces the trivial operation on $\unsure{\mathcal G}$ --- which is indeed trivial by (3) a).

We can therefore set
\[
(6)\qquad  G\overset{\unsure{\mathrm{d\acute ef}|\mathrm{dfn}}}{=}S_0G/i(SN),
\]
we then find a commutative diagram

\begin{samepage}
\noindent (7)\quad
\diagram{\begin{tikzcd}[column sep=small,row sep=small]
& N \arrow[dr,"\varphi_N"] & \\
SN \arrow[ur,hook] \arrow[dr,"\varphi_{SN}"'] & & G \\
& \unsure{\mathcal G} \arrow[ur,hook] &
\end{tikzcd}}
\end{samepage},

we see that the image of $\unsure{\mathcal G}$ in $G$ is an invariant subgroup, and that \unsure{we} have an isomorphism on the quotients:
\[
(8)\qquad  N/SN\xrightarrow{\sim}G/\unsure{\mathcal G},
\]
which implies that we also have an isomorphism
\[
(9)\qquad  \operatorname{Ker}(\varphi_{SN}:SN\longrightarrow\unsure{\mathcal G})
\xrightarrow{\sim}\operatorname{Ker}(N\xrightarrow{\varphi_N}G);
\]
\deleted{thus $\varphi_N$}

\begin{samepage}
\noindent (10)\quad
\diagram{\(\begin{array}{l}
\varphi_N:N\longrightarrow G\text{ is an embedding}\\
\unsure{\Updownarrow}\\
\varphi_{SN}:SN\longrightarrow\unsure{\mathcal G}\text{ is an embedding.}
\end{array}\)}
\end{samepage}

We thus find \uline{equivalences of categories}, inverse to one another, between \inserted{\unsure{the categories}} of

\margin{692}

% Printed archive footer.
\begin{center}
692\\
\texttt{https://grothendieck.umontpellier.fr}\\
Université de Montpellier, UMR CNRS 5149, Case courrier 051, Place Eugène Bataillon\\
34095 Montpellier cedex 5 - France
\end{center}


\scanpage{694}

situations (1), and that of situations (2).

\textcircled{b}
Let us now consider a situation
\[
(11)\qquad 
\left\{\begin{array}{l}
G,\quad \unsure{\mathcal G}\subset G,\quad (\varphi_i:N_i\longrightarrow G)_{i\in I}\\[4pt]
\text{hence }SN_i=\varphi_i^{-1}(\unsure{\mathcal G})
\end{array}\right.
\]
under the condition
\[
(12)\qquad 
\forall i\in I,\quad \varphi_i:N_i\longrightarrow G/\unsure{\mathcal G}=\gamma\quad\text{epimorphism}
\]
\unsure{can} we reconstruct it from
\[
(13)\qquad 
\unsure{\mathcal G},\ (N_i,SN_i,\unsure{\varphi_i},\Theta_i)_{i\in I}
\]
where \inserted{$\forall i\in I,$} $(N_i,SN_i,\varphi_i,\Theta_i)$ satisfies the conditions (3) a) and b) --- \unsure{where|and} of course $\forall i\in I$
\[
(14)\qquad 
\deleted{\unsure{\varphi_i}}\inserted{\Theta_i}:N_i\longrightarrow\operatorname{Aut}(\unsure{\mathcal G}),\qquad
\varphi_i:SN_i\longrightarrow\unsure{\mathcal G}.
\]

It is tempting to reconstruct $G$ with the help \unsure{of|of the} $\unsure{\mathcal G}$ and of the $N_i$, as a sort of ``amalgamated sum'' relative to:

\noindent \deleted{\unsure{(15)}}\quad
\diagram{\begin{tikzcd}[column sep=small,row sep=small]
& \unsure{\mathcal G} \arrow[dr,dashed] & \\
SN_i \arrow[ur,"\varphi_i"] \arrow[dr,hook,"\mathrm{plongt}"'] & & G \\
& N_i \arrow[ur,dashed] &
\end{tikzcd}}

\margin{693}

% Printed archive footer.
\begin{center}
693\\
\texttt{https://grothendieck.umontpellier.fr}\\
Université de Montpellier, UMR CNRS 5149, Case courrier 051, Place Eugène Bataillon\\
34095 Montpellier cedex 5 - France
\end{center}


\scanpage{695}
\mspage{650}

\deleted{\uline{Proposition} \unsure{Starting|Start} from (11), hence (13)}

% small x-shaped marks follow "dire," and $\Theta_i$ in this sentence (reference marks?); not transcribed as text.
But it is more plausible to say, that for every $i$, we reconstruct \unsure{via|here} $N_i,SN_i,\varphi_i,\Theta_i$ the situation (for $i$ fixed)

\noindent \deleted{\unsure{(17)}}\inserted{\unsure{(15)}}\quad
\diagram{\begin{tikzcd}
\deleted{\unsure{\varphi_i}}N_i \arrow[r] & G_i(\unsure{\simeq}G) \\
SN_i \arrow[u,hook] \arrow[r] & \unsure{\mathcal G} \arrow[u,hook]
\end{tikzcd}}

and that we reconstruct (11) itself by giving ourselves a \uline{transitive system} \uline{of isomorphisms} \uline{between the} $G_i$ $(i\in I)$, \uline{inducing} \uline{the identity on} $\unsure{\mathcal G}$. Such a datum implies in the first place the datum of a transitive system of isomorphisms between the
\[
(16)\qquad 
\gamma_i\overset{\unsure{\mathrm{d\acute ef}|\mathrm{dfn}}}{=}N_i/SN_i
\]
--- which we \unsure{will identify} therefore with one and the same group
\[
(17)\qquad 
\begin{gathered}
\gamma,\text{ common quotient group}\\
\text{of the }N_i\text{ (by the }SN_i\text{).}
\end{gathered}
\]
Let us now identify the $G/\unsure{\mathcal G}_i=\gamma_i$ with $\gamma$
\[
(18)\qquad G_i/\unsure{\mathcal G}\simeq\gamma,
\]
and form the fiber product

\margin{694}

% Printed archive footer.
\begin{center}
694\\
\texttt{https://grothendieck.umontpellier.fr}\\
Université de Montpellier, UMR CNRS 5149, Case courrier 051, Place Eugène Bataillon\\
34095 Montpellier cedex 5 - France
\end{center}


\scanpage{696}

\[
(19)\qquad 
\unsure{g}=\underset{\text{(fiber product over }\gamma\text{)}}{\prod_{i\in I}\!{}_{\gamma}G_i}\ \supset\ \unsure{\mathcal G}^I
\]
\margin{$\unsure{g}$ is an extension of $\gamma$ \unsure{by} $\unsure{\mathcal G}^I$}

The datum of a transitive system of isomorphisms between the $G_i$, inducing the identity on the subgroups $\unsure{\mathcal G}$, and compatible with the system \deleted{\unsure{of isomorphisms|of the}} \unsure{transitive} of isomorphisms already considered between the $\gamma_i=N_i/SN_i\simeq G_i/\unsure{\mathcal G}$, is then equivalent to the datum of a subgroup $G$ of $\unsure{g}$, satisfying the conditions
\[
(20)\qquad 
\left\{\begin{array}{ll}
\text{a)}&G\cap\unsure{\mathcal G}^I=\text{diagonal }\delta(\unsure{\mathcal G})\text{ of }\unsure{\mathcal G}^I\\[5pt]
\text{b)}&\text{The projection }G\longrightarrow\gamma\text{ is an epimorphism}
\end{array}\right.
\]

In other words, we have constructed, with the help of (13) and of the transitive system of isomorphisms between the $\gamma_i$, an extension $\unsure{g}$ of their ``common value'' $\gamma$, by $\unsure{\mathcal G}^I$ --- and the additional structure, corresponding to the datum \unsure{(11)}, is the ``restriction'' of this extension to an extension of $\gamma$ by the diagonal subgroup of $\unsure{\mathcal G}^I$.

\margin{695}

% Printed archive footer.
\begin{center}
695\\
\texttt{https://grothendieck.umontpellier.fr}\\
Université de Montpellier, UMR CNRS 5149, Case courrier 051, Place Eugène Bataillon\\
34095 Montpellier cedex 5 - France
\end{center}
